\textbf{Working methods}: The exam text, the test files of the exercises
and a synoptic table on Python syntax can be found in the \texttt{exam}
folder on your desk, in the ``Statistica-Laboratorio'' virtual machine.
The solutions to the exercises must be placed in the same folder
\texttt{exam}. \textbf{First compile the file} \texttt{studente.txt}
with the first name, surname and matriculation information

\begin{verbatim}
nome: Giacomo
cognome: Leopardi
matricola: 123456789
\end{verbatim}

Not completing the \texttt{student.txt} file is like not handing in the
assignment. For the exam, you can only have

\begin{itemize}
\tightlist
\item
  your ID or passport;
\item
  beverages and snacks;
\end{itemize}

and you must leave your bags and backpacks in the corner of the room,
with mobile phones and electronic devices switched off inside.

\textbf{The grade is the sum of the points of the exercises}, so it is
possible to get 30 even without completing them all.

\textbf{Use of test files}: tests check very briefly the setting of the
solutions and passing the tests does not at all imply that the solution
is correct. Read the text of the exercises before starting to solve
them. To complete each exercise exercise you must

\begin{itemize}
\tightlist
\item
  write a solution file \textbf{with the specified name};
\item
  have the function \textbf{with the specified name} inside;
\item
  place the solution file in the \texttt{exam} folder;
\item
  run the \texttt{python3\ \textless{}testfile\textgreater{}} command in
  that folder
\end{itemize}

where \texttt{\textless{}testfile\textgreater{}} should of course be
replaced with the name of the appropriate test file for the exercise you
are working on.

There is an independent test file for each exercise, so you can work on
the exercises one at a time with more ease. The result of each test is a
screen (or several screens) showing

\begin{itemize}
\tightlist
\item
  whether the file with the correct name was found,
\item
  whether the file contains the function with the correct name,
\item
  whether calls to the function with different parameter values end by
  returning correct results.
\end{itemize}

For each function written, calls are made with different values of the
parameters. The result of the tests is returned with the character

\begin{itemize}
\tightlist
\item
  \texttt{E} if the call cannot be executed,
\item
  \texttt{F} if the function does not return the correct result,
\item
  \texttt{.} if the function returns the correct result.
\end{itemize}

\pagebreak

\section{Circle area}\label{circle-area}

Write a function \texttt{circle\_area(x)} that:

\begin{itemize}
\tightlist
\item
  takes a non-negative number \texttt{x} as parameter;
\item
  returns the area of a circle with radius \texttt{x}.
\end{itemize}

After importing \texttt{math} package, a good approximation of Pi is
available as \texttt{math.pi}

The function must be saved in a file named \texttt{circle\_area.py}

The Python function must be saved in file: \texttt{circle\_area.py}
\Test file: \texttt{test\_circle\_area.py}

\pagebreak

\section{Average}\label{average}

Write a function \texttt{average(L)} that,

\begin{itemize}
\tightlist
\item
  takes a list \texttt{L} as parameter, in which each element is a
  number (\texttt{int} or \texttt{float});
\item
  returns the average of asll numbers in the list \texttt{L}. You can
  assume that \texttt{L} is a list of numbers (each element is
  \texttt{int} oppure \texttt{float}).
\end{itemize}

The function must be saved in a file named \texttt{average.py}

The Python function must be saved in file: \texttt{average.py}
\Test file: \texttt{test\_average.py}

\pagebreak

\section{Scalar product between two
vectors}\label{scalar-product-between-two-vectors}

The scalar product between two vectors
\(\langle a_0, a_1,\ldots, a_{n-1}  \rangle\) and
\(\langle b_0,b_1,\ldots, b_{n-1} \rangle\) is defined as

\begin{equation*}
      \sum_{i=0}^{n-1} a_i b_i
\end{equation*}

We represent vectors as lists of numbers.

Write a function \texttt{scalar\_product} that:

\begin{itemize}
\item
  takes two non-empty lists \texttt{a}, \texttt{b} with the same length
  as parameters, where each element is a number (\texttt{int} or
  \texttt{float});
\item
  returns the scalar product of the two lists.
\end{itemize}

The scalar product between two lists is therefore the sum of the
products of the homologous elements. For example, the scalar product
between the vectors \texttt{{[}1,\ 4,\ 3{]}} and
\texttt{{[}5,\ 5,\ 1{]}} is
\(1 \times 5 + 4 \times 5 + 3 \times 1 = 28\). The scalar product of two
vectors is a numerical value.

You can \textbf{assume} that the lists are not empty, that they contain
only numbers and that they have the same length.

The function must be saved in a file named \texttt{scalar\_product.py}

The Python function must be saved in file: \texttt{scalar\_product.py}
\Test file: \texttt{test\_scalar\_product.py}

\pagebreak

\section{Check that all values in a dictionary are
small}\label{check-that-all-values-in-a-dictionary-are-small}

Write a function \texttt{dict\_check\_small\_values(d,\ threshold)}
that:

\begin{itemize}
\tightlist
\item
  takes as arguments a dictionary \texttt{d} with values of type
  \texttt{int} or \texttt{float} and an number \texttt{threshold};
\item
  returns \texttt{True} if all values are strictly smaller than
  \texttt{threshold}, and returns \texttt{False} otherwise.
\end{itemize}

You may assume that alle values are numeric.

The function must be saved in a file named \texttt{scalar\_product}.

For example:

\begin{Shaded}
\begin{Highlighting}[]
\NormalTok{d }\OperatorTok{=}\NormalTok{ \{}\StringTok{"Mario"}\NormalTok{: }\DecValTok{1000}\NormalTok{, }\StringTok{"Gianna"}\NormalTok{ : }\DecValTok{8000}\NormalTok{, }\StringTok{"Alberto"}\NormalTok{: }\DecValTok{4000}\NormalTok{\}}

\BuiltInTok{print}\NormalTok{(dict\_check\_small\_values(d, }\DecValTok{7000}\NormalTok{))}
\end{Highlighting}
\end{Shaded}

\begin{verbatim}
False
\end{verbatim}

The Python function must be saved in file:
\texttt{dict\_check\_small\_values.py} \Test file:
\texttt{test\_dict\_check\_small\_values.py}

\pagebreak

\section{Count large elements in a
matrix}\label{count-large-elements-in-a-matrix}

Write a function \texttt{mat\_count\_large(m,\ threshold)} that:

\begin{itemize}
\tightlist
\item
  takes as argument a matrix \texttt{m}, represented as a list of lists;
\item
  returns the number of elements in \texttt{m} that are strictly larger
  than \texttt{threshold}.
\end{itemize}

We represent an \(r \times c\) matrix in Python as a list with length
\(r\), where each element is a list of length \(c\). You may assume that
the argument is well formed, with \(r \geq 1\), \(c \geq 1\), and all
the lists representing rows contain \texttt{c}~numeric elements.

The function must be saved in a file named \texttt{mat\_count\_large}.

The Python function must be saved in file: \texttt{mat\_count\_large.py}
\Test file: \texttt{test\_mat\_count\_large.py}

\pagebreak

\section{Change minima to zero}\label{change-minima-to-zero}

Write a function \texttt{change\_mins\_in\_zero(values)} that:

\begin{itemize}
\tightlist
\item
  expects a non-empty list of comparable values as an argument;
\item
  directly modifies the \texttt{values} list by replacing the minimum
  values with zero.
\end{itemize}

If the minimum appears several times in the list, the function must
replace all occurrences of the minimum value with zero.

The function must be saved in a file named \texttt{mat\_count\_large}.

The Python function must be saved in file:
\texttt{change\_mins\_in\_zero.py} \Test file:
\texttt{test\_change\_mins\_in\_zero.py}
